% markpage — export LaTeX (cible : xelatex ou lualatex).
% La compilation avec pdflatex échouera : fontspec + UTF-8 natif
% dans les blocs de code requièrent xelatex / lualatex.
\documentclass[11pt,a4paper]{article}
\usepackage{fontspec}
\IfFontExistsTF{DejaVu Serif}{\setmainfont{DejaVu Serif}}{}
\IfFontExistsTF{DejaVu Sans}{\setsansfont{DejaVu Sans}}{}
\IfFontExistsTF{DejaVu Sans Mono}{\setmonofont{DejaVu Sans Mono}}{}

\IfFileExists{french.ldf}{\usepackage[french]{babel}}{}
\usepackage{amsmath,amssymb,amsthm}
\usepackage{stmaryrd}
\usepackage{graphicx}
\usepackage[export]{adjustbox}
\usepackage{hyperref}
\usepackage{xcolor}
\usepackage{listings}
\usepackage{enumitem}
\usepackage{booktabs}
\usepackage{tabularx}
\usepackage[normalem]{ulem}
\usepackage[breakable,skins]{tcolorbox}
\usepackage{newunicodechar}

% DejaVu Serif (our main text font) doesn't ship the Unicode
% math-bracket / logic glyphs that show up in our prose (e.g.
% explanations of the editor's ligatures). Redirect each known
% glyph to its LaTeX command via \ensuremath, so it renders
% from the math font even when typed in normal text.
\newunicodechar{⟦}{\ensuremath{\llbracket}}
\newunicodechar{⟧}{\ensuremath{\rrbracket}}
\newunicodechar{⟨}{\ensuremath{\langle}}
\newunicodechar{⟩}{\ensuremath{\rangle}}
\newunicodechar{⊢}{\ensuremath{\vdash}}
\newunicodechar{⊣}{\ensuremath{\dashv}}
\newunicodechar{⊨}{\ensuremath{\models}}
\newunicodechar{⊥}{\ensuremath{\bot}}
\newunicodechar{⊤}{\ensuremath{\top}}
\newunicodechar{∀}{\ensuremath{\forall}}
\newunicodechar{∃}{\ensuremath{\exists}}
\newunicodechar{∈}{\ensuremath{\in}}
\newunicodechar{∉}{\ensuremath{\notin}}
\newunicodechar{∅}{\ensuremath{\emptyset}}
\newunicodechar{⊂}{\ensuremath{\subset}}
\newunicodechar{⊆}{\ensuremath{\subseteq}}
\newunicodechar{⊃}{\ensuremath{\supset}}
\newunicodechar{⊇}{\ensuremath{\supseteq}}
\newunicodechar{∪}{\ensuremath{\cup}}
\newunicodechar{∩}{\ensuremath{\cap}}
\newunicodechar{∧}{\ensuremath{\land}}
\newunicodechar{∨}{\ensuremath{\lor}}
\newunicodechar{¬}{\ensuremath{\neg}}
\newunicodechar{★}{\ensuremath{\bigstar}}
\newunicodechar{✓}{\ensuremath{\checkmark}}
\newunicodechar{✗}{\ensuremath{\times}}
\newunicodechar{♥}{\ensuremath{\heartsuit}}
\newunicodechar{♦}{\ensuremath{\diamondsuit}}
\newunicodechar{♠}{\ensuremath{\spadesuit}}
\newunicodechar{♣}{\ensuremath{\clubsuit}}
\newunicodechar{⚠}{\textbf{[!]}}

\hypersetup{colorlinks=true, linkcolor=blue!50!black, urlcolor=blue!50!black}
\lstset{
  basicstyle=\ttfamily\small,
  breaklines=true,
  frame=single,
  framesep=4pt,
}

% Theorem-like environments matching the ::: admonitions (§17 of
% the markpage SPEC). The shared counter `theorem` keeps lemma /
% proposition / corollary numbered in the same sequence as the
% theorems themselves, which is the usual mathematical convention.
\newtheorem{theorem}{Théorème}
\newtheorem{lemma}[theorem]{Lemme}
\newtheorem{proposition}[theorem]{Proposition}
\newtheorem{corollary}[theorem]{Corollaire}
\theoremstyle{definition}
\newtheorem{definition}{Définition}
\newtheorem{example}{Exemple}
\newtheorem{remark}{Remarque}

\title{Cross-references \textbackslash{}textbackslash\{\}label\textbackslash{}\{sec:intro\textbackslash{}\}}
\author{Test Author \\ Test Org}
\date{2026-01-01}

\begin{document}
\maketitle

Cross-references follow LaTeX conventions: a \texttt{\textbackslash{}label\{key\}} attaches an
anchor to a numbered target, \texttt{\textbackslash{}ref\{key\}} resolves to the number with a
clickable link. The user writes the prefix word themselves so the
grammar stays natural — "voir algorithme \textbackslash{}ref\{alg:bubble\}" rather
than a frozen "Algorithme 1".


\section{Captioned blocks \textbackslash{}label\{sec:blocks\}}

A figure carries a label after the caption:


\begin{lstlisting}
S
  NP
    Det
    N
  VP
    V
    NP
      Det
      N
\end{lstlisting}


A table:


\begin{center}
\begin{tabular}{ll}
  \toprule
  Équipe & Effectif \\
  \midrule
  Backend & 12 \\
  Frontend & 8 \\
  \bottomrule
\end{tabular}
\end{center}


An algorithm:


\begin{lstlisting}
for i from 1 to n - 1 do
  for j from 0 to n - i - 1 do
    if A[j] > A[j + 1] then
      swap A[j] and A[j + 1]
    end
  end
end
return A
\end{lstlisting}


A code listing:


\begin{lstlisting}[language=Python]
def median(xs):
    s = sorted(xs)
    n = len(s)
    return s[n // 2] if n % 2 else (s[n // 2 - 1] + s[n // 2]) / 2
\end{lstlisting}


\section{Équations \textbackslash{}label\{sec:equations\}}

Une équation étiquetée se voit attribuer un numéro à droite :


\[
\int_0^1 x^2 \, dx = \frac{1}{3} \label{eq:cube}
\]


Tandis qu'une équation non étiquetée n'en reçoit pas :


\[
e^{i\pi} + 1 = 0
\]


Une seconde équation étiquetée prend le numéro suivant dans la
séquence des équations numérotées :


\[
\sum_{k=1}^n k = \frac{n(n+1)}{2} \label{eq:gauss}
\]


\section{Références croisées \textbackslash{}label\{sec:refs\}}

On peut maintenant écrire : « \textbackslash{}ref\{sec:intro\} » introduit la
notation, la figure \textbackslash{}ref\{fig:syntax\} illustre la décomposition,
le tableau \textbackslash{}ref\{tab:teams\} donne les effectifs, l'algorithme
\textbackslash{}ref\{alg:bubble\} décrit le tri, et le listing \textbackslash{}ref\{lst:median\}
calcule la médiane. Voir aussi l'équation \textbackslash{}ref\{eq:cube\} pour
l'intégrale, et l'équation \textbackslash{}ref\{eq:gauss\} pour la somme. Les
références aux \textbf{sections} rendent le titre lui-même (pratique
quand les sections ne sont pas numérotées), tandis que celles aux
\textbf{figures / tableaux / algos / listings / équations} rendent le
numéro.


\section{Erreur de référence}

Une référence inconnue (typo) est rendue visible en rouge avec
un tooltip : voir la section \textbackslash{}ref\{sec:typo-inexistante\}.

\end{document}
